% !TEX root = ../DoAn.tex
\documentclass[../DoAn.tex]{subfiles}
\begin{document}

Chương này trình bày bảy nền tảng lý thuyết trực tiếp định hình các quyết định thiết kế
then chốt của hệ thống. Mỗi mục xuất phát từ một bài toán kỹ thuật cụ thể mà hệ thống
phải giải quyết, sau đó phân tích nguyên lý và đánh giá cách áp dụng: (1)~kiến trúc
microservices là nền tảng triển khai toàn bộ hệ thống; (2)~phân rã hệ thống và cách ly
dữ liệu dẫn đến bài toán nhất quán phân tán; (3)~giao dịch nhiều bước xuyên service đòi
hỏi Saga Orchestration; (4)~giao tiếp bất đồng bộ tin cậy cần Transactional Outbox và
Idempotent Consumer; (5)~flash sale đồng thời cao yêu cầu thao tác nguyên tử trên Redis;
(6)~xác thực phân tán cần mô hình Trusted Subsystem; (7)~các mẫu hỗ trợ bổ sung hoàn
thiện độ tin cậy và khả năng quan sát.

% =======================================================================
% \section{Kiến trúc Microservices}
% \label{section:2.1}

% \subsection{Khái niệm và đặc điểm cốt lõi}

% Kiến trúc microservices là phong cách kiến trúc phần mềm cấu trúc một ứng dụng thành tập
% hợp các service nhỏ, triển khai độc lập, mỗi service chạy trong tiến trình riêng và giao
% tiếp qua cơ chế nhẹ --- thường là HTTP/REST hoặc message broker. Mỗi service được tổ chức
% xung quanh một khả năng nghiệp vụ (business capability) cụ thể và có thể được triển khai
% độc lập bằng hệ thống tự động hóa~\cite{fowler2014, newman2021}.

% Sáu đặc điểm cốt lõi phân biệt microservices với các kiểu kiến trúc khác:

% \begin{itemize}[noitemsep]
%     \item \textbf{Single Responsibility}: mỗi service đảm nhiệm đúng một bounded context
%           nghiệp vụ, không chồng lấn trách nhiệm. Ranh giới rõ ràng giúp nhóm phát triển
%           hiểu toàn bộ service và thay đổi mà không ảnh hưởng service khác.
%     \item \textbf{Loose Coupling, High Cohesion}: các service không biết chi tiết nội bộ
%           của nhau (loose coupling), mọi giao tiếp đều qua API/sự kiện được công bố công
%           khai; trong khi đó logic liên quan được nhóm lại trong cùng service (high cohesion).
%     \item \textbf{Decentralized Data Management}: mỗi service sở hữu và quản lý kho dữ
%           liệu riêng (Database per Service). Không có service nào được phép đọc trực tiếp
%           vào database của service khác.
%     \item \textbf{Independent Deployability}: mỗi service có pipeline CI/CD riêng, có thể
%           cập nhật, rollback và scale mà không cần phối hợp với service khác.
%     \item \textbf{Technology Heterogeneity}: mỗi service được phép chọn ngôn ngữ lập trình,
%           framework và loại cơ sở dữ liệu phù hợp nhất với đặc thù nghiệp vụ của nó
%           (Polyglot Persistence).
%     \item \textbf{Design for Failure}: trong hệ thống phân tán, sự cố mạng, timeout và
%           lỗi từng phần là tất yếu. Mỗi service phải được thiết kế giả định service lân
%           cận có thể không sẵn sàng và xử lý gracefully qua Circuit Breaker, Retry và
%           fallback.
% \end{itemize}

% \subsection{So sánh với Monolith và SOA}

% Để đánh giá đúng lý do lựa chọn microservices, cần đối chiếu với hai kiểu kiến trúc tiền
% nhiệm là Monolith và SOA (Service-Oriented Architecture). Bảng~\ref{table:arch-comparison}
% tổng hợp sự khác biệt theo các chiều quan trọng:

% \begin{table}[H]
% \centering
% \caption{So sánh kiến trúc Monolith, SOA và Microservices}
% \label{table:arch-comparison}
% \begin{compacttable}
% \begin{tabularx}{\textwidth}{|L{3.2cm}|Y|Y|Y|}
% \hline
% \textbf{Tiêu chí} & \textbf{Monolith} & \textbf{SOA} & \textbf{Microservices} \\ \hline
% Đơn vị triển khai
%     & Toàn bộ ứng dụng trong một artifact duy nhất
%     & Nhóm chức năng, thường là WAR/EAR lớn
%     & Service đơn lẻ, JAR độc lập, container riêng \\ \hline
% Cơ sở dữ liệu
%     & Shared database duy nhất
%     & Shared hoặc ít database dùng chung
%     & Database per Service --- mỗi service sở hữu DB riêng \\ \hline
% Giao tiếp nội bộ
%     & Gọi hàm trong bộ nhớ (in-process)
%     & SOAP/XML qua Enterprise Service Bus (ESB)
%     & REST, gRPC hoặc message broker nhẹ (Kafka) \\ \hline
% Công nghệ
%     & Bị ràng buộc một stack duy nhất
%     & Chủ yếu một ngôn ngữ, ESB là nút thắt
%     & Tự do --- polyglot theo từng service \\ \hline
% Khả năng mở rộng
%     & Scale toàn bộ ứng dụng, lãng phí tài nguyên
%     & Scale theo nhóm service
%     & Scale từng service độc lập theo nhu cầu \\ \hline
% Tổ chức nhóm
%     & Nhóm lớn, chung codebase, dễ xung đột
%     & Nhóm trung bình theo layer
%     & Nhóm nhỏ, mỗi nhóm sở hữu vài service \\ \hline
% Độ phức tạp vận hành
%     & Thấp --- một artifact, một deployment
%     & Trung bình --- ESB là điểm phức tạp
%     & Cao --- phân tán, cần observability, service mesh \\ \hline
% Phù hợp với
%     & MVP, nhóm nhỏ, ứng dụng đơn giản
%     & Enterprise legacy, tích hợp hệ thống cũ
%     & Hệ thống lớn, nhiều nhóm, yêu cầu scale cao \\ \hline
% \end{tabularx}
% \end{compacttable}
% \end{table}

% \textbf{Kiến trúc Monolith} đơn giản trong triển khai và debug --- toàn bộ code chạy
% trong cùng một tiến trình nên không có độ trễ mạng nội bộ, không cần service discovery
% và transaction ACID hoạt động tự nhiên. Tuy nhiên, khi hệ thống tăng trưởng, monolith
% gặp hai vấn đề không thể tránh: mọi thay đổi dù nhỏ đều yêu cầu rebuild và redeploy
% toàn bộ ứng dụng; scale chỉ có thể thực hiện theo chiều ngang toàn bộ ứng dụng, không
% thể scale riêng module nào đang là bottleneck.

% \textbf{SOA} giải quyết một phần bằng cách chia nhỏ thành các service lớn hơn, nhưng
% phụ thuộc nặng vào ESB như một bus trung tâm. ESB trở thành điểm tập trung cho routing,
% transformation và orchestration --- khi ESB gặp sự cố hoặc trở thành bottleneck, toàn bộ
% hệ thống bị ảnh hưởng. Ngoài ra, shared database trong SOA vẫn tạo ra coupling ở tầng dữ
% liệu ngay cả khi logic đã được tách ra.

% \textbf{Microservices} đẩy nguyên tắc tách biệt đến mức tối đa: mỗi service là một đơn
% vị triển khai độc lập hoàn toàn, từ code đến database. Đánh đổi là độ phức tạp vận hành
% tăng đáng kể --- phải giải quyết các bài toán phân tán như nhất quán dữ liệu, distributed
% tracing và service-to-service authentication mà monolith không bao giờ gặp phải. Các
% phần tiếp theo của chương phân tích cụ thể từng bài toán đó.

% Hệ thống được xây dựng theo kiến trúc microservices với 13 business service và 3 service
% hạ tầng. Mỗi business service tương ứng với một domain nghiệp vụ tách biệt: định danh
% (identity-service), hồ sơ người dùng (user-service), catalog sản phẩm (product-service),
% tìm kiếm (search-service), giỏ hàng (cart-service), đặt hàng (order-service), thanh toán
% (payment-service), kho hàng (inventory-service), voucher (voucher-service), đánh giá
% (review-service), nội dung (content-service), thông báo (notification-service) và flash
% sale (flash-sale-service). Ba service hạ tầng --- API Gateway, Config Server và Service
% Registry --- hỗ trợ vận hành mà không chứa logic nghiệp vụ. Mỗi
% service được đóng gói thành Docker container riêng với cơ sở dữ liệu độc lập, pipeline
% CI/CD riêng và giao tiếp với nhau qua REST (đồng bộ) hoặc Kafka (bất đồng bộ). Quyết định
% này dẫn trực tiếp đến các thách thức kỹ thuật được giải quyết trong các phần tiếp theo.

\section{Kiến trúc Microservices}
\label{section:2.1}

\subsection{Khái niệm và đặc điểm cốt lõi}

Kiến trúc microservices là phong cách kiến trúc phần mềm cấu trúc một ứng dụng thành tập
hợp các service nhỏ, triển khai độc lập, mỗi service chạy trong tiến trình riêng và giao
tiếp qua cơ chế nhẹ --- thường là HTTP/REST hoặc message broker. Mỗi service được tổ chức
xung quanh một khả năng nghiệp vụ (business capability) cụ thể và có thể được triển khai
độc lập bằng hệ thống tự động hóa~\cite{fowler2014, newman2021}.

Sáu đặc điểm cốt lõi phân biệt microservices với các kiểu kiến trúc khác:

\begin{itemize}[noitemsep]
    \item \textbf{Single Responsibility}: mỗi service đảm nhiệm đúng một bounded context
          nghiệp vụ, không chồng lấn trách nhiệm. Ranh giới rõ ràng giúp nhóm phát triển
          hiểu toàn bộ service và thay đổi mà không ảnh hưởng service khác.
    \item \textbf{Loose Coupling, High Cohesion}: các service không biết chi tiết nội bộ
          của nhau (loose coupling), mọi giao tiếp đều qua API hoặc sự kiện được công bố
          công khai; trong khi đó logic liên quan được nhóm lại trong cùng service
          (high cohesion).
    \item \textbf{Decentralized Data Management}: mỗi service sở hữu và quản lý kho dữ
          liệu riêng (Database per Service). Không service nào được phép đọc trực tiếp vào
          database của service khác.
    \item \textbf{Independent Deployability}: mỗi service có pipeline CI/CD riêng, có thể
          cập nhật, rollback và scale mà không cần phối hợp với service khác.
    \item \textbf{Technology Heterogeneity}: mỗi service được phép chọn ngôn ngữ lập trình,
          framework và loại cơ sở dữ liệu phù hợp nhất với đặc thù nghiệp vụ của nó
          (Polyglot Persistence).
    \item \textbf{Design for Failure}: trong hệ thống phân tán, sự cố mạng, timeout và
          lỗi từng phần là tất yếu. Mỗi service phải được thiết kế với giả định rằng
          service lân cận có thể không sẵn sàng, và xử lý tình huống đó một cách an toàn
          (graceful) thông qua các cơ chế như Circuit Breaker, Retry và fallback.
\end{itemize}

\subsection{So sánh với Monolith và SOA}

Để đánh giá đúng lý do lựa chọn microservices, cần đối chiếu với hai kiểu kiến trúc tiền
nhiệm là Monolith và SOA (Service-Oriented Architecture). Bảng~\ref{table:arch-comparison}
tổng hợp sự khác biệt theo các chiều quan trọng:

\begin{table}[H]
\centering
\caption{So sánh kiến trúc Monolith, SOA và Microservices}
\label{table:arch-comparison}
\begin{compacttable}
\begin{tabularx}{\textwidth}{|L{3.2cm}|Y|Y|Y|}
\hline
\textbf{Tiêu chí} & \textbf{Monolith} & \textbf{SOA} & \textbf{Microservices} \\ \hline
Đơn vị triển khai
    & Toàn bộ ứng dụng trong một artifact duy nhất
    & Nhóm chức năng, thường là WAR/EAR lớn
    & Service đơn lẻ, JAR độc lập, container riêng \\ \hline
Cơ sở dữ liệu
    & Shared database duy nhất
    & Shared hoặc ít database dùng chung
    & Database per Service --- mỗi service sở hữu DB riêng \\ \hline
Giao tiếp nội bộ
    & Gọi hàm trong bộ nhớ (in-process)
    & SOAP/XML qua Enterprise Service Bus (ESB)
    & REST, gRPC hoặc message broker nhẹ (Kafka) \\ \hline
Công nghệ
    & Bị ràng buộc một stack duy nhất
    & Chủ yếu một ngôn ngữ, ESB là nút thắt
    & Tự do --- polyglot theo từng service \\ \hline
Khả năng mở rộng
    & Scale toàn bộ ứng dụng, lãng phí tài nguyên
    & Scale theo nhóm service
    & Scale từng service độc lập theo nhu cầu \\ \hline
Tổ chức nhóm
    & Nhóm lớn, chung codebase, dễ xung đột
    & Nhóm trung bình theo layer
    & Nhóm nhỏ, mỗi nhóm sở hữu vài service \\ \hline
Độ phức tạp vận hành
    & Thấp --- một artifact, một deployment
    & Trung bình --- ESB là điểm phức tạp
    & Cao --- phân tán, cần observability, service mesh \\ \hline
Phù hợp với
    & MVP, nhóm nhỏ, ứng dụng đơn giản
    & Enterprise legacy, tích hợp hệ thống cũ
    & Hệ thống lớn, nhiều nhóm, yêu cầu scale cao \\ \hline
\end{tabularx}
\end{compacttable}
\end{table}

\textbf{Kiến trúc Monolith} đơn giản trong triển khai và gỡ lỗi --- toàn bộ code chạy
trong cùng một tiến trình nên không có độ trễ mạng nội bộ, không cần service discovery và
transaction ACID hoạt động tự nhiên. Tuy nhiên, khi hệ thống tăng trưởng, monolith gặp hai
vấn đề khó tránh: mọi thay đổi dù nhỏ đều yêu cầu rebuild và redeploy toàn bộ ứng dụng; và
việc mở rộng chỉ có thể thực hiện cho toàn bộ ứng dụng, không thể scale riêng module đang
là bottleneck.

\textbf{SOA} giải quyết một phần bằng cách chia hệ thống thành các service lớn hơn, nhưng
phụ thuộc nặng vào ESB như một bus trung tâm. ESB đảm nhiệm routing, transformation và
orchestration --- do đó khi ESB gặp sự cố hoặc trở thành bottleneck, toàn bộ hệ thống bị
ảnh hưởng. Ngoài ra, việc dùng chung database trong SOA vẫn tạo ra coupling ở tầng dữ liệu
ngay cả khi logic đã được tách rời.

\textbf{Microservices} đẩy nguyên tắc tách biệt đến mức tối đa: mỗi service là một đơn vị
triển khai độc lập hoàn toàn, từ code đến database. Đánh đổi cho sự độc lập đó là độ phức
tạp vận hành tăng đáng kể --- hệ thống phải giải quyết các bài toán phân tán như nhất quán
dữ liệu, distributed tracing và xác thực giữa các service, vốn không tồn tại trong monolith.
Các phần tiếp theo của chương sẽ lần lượt phân tích những bài toán này cùng các mẫu thiết
kế tương ứng.

Các nguyên tắc trên là nền tảng định hướng cho hệ thống của đề tài: hệ thống được phân rã thành các business service theo từng miền nghiệp vụ, kết hợp với một số service hạ tầng, trong đó mỗi service sở hữu cơ sở dữ liệu riêng và giao tiếp qua REST hoặc Kafka.
% =======================================================================
% \section{Phân rã hệ thống và Database per Service}
% \label{section:2.2}

% Hệ thống thương mại điện tử phức tạp không thể được mô hình hóa bằng một mô hình dữ liệu
% thống nhất duy nhất, vì các miền nghiệp vụ khác nhau nhìn nhận cùng một khái niệm theo những
% cách không tương thích. Ví dụ, khái niệm \textit{Đơn hàng} trong miền đặt hàng mang ý nghĩa
% về thông tin khách hàng, danh sách sản phẩm và phương thức thanh toán; trong miền kho hàng,
% chính đơn hàng đó lại là tập hợp các mặt hàng cần trừ tồn kho. Domain-Driven Design (DDD)
% giải quyết mâu thuẫn này thông qua khái niệm \textbf{Bounded Context}~\cite{evans2003}: một ranh giới tường
% minh trong đó một mô hình nghiệp vụ cụ thể có ý nghĩa nhất quán và không cần giải thích thêm.

% Áp dụng DDD vào hệ thống, 13 Bounded Context được xác định tương ứng với 13 microservice:
% mỗi service sở hữu mô hình nghiệp vụ riêng, không chia sẻ lớp domain với service khác.
% Ranh giới giữa các context được thể hiện qua nguyên tắc \textit{Database per Service}~\cite{richardson2018}: mỗi
% service chỉ được phép đọc và ghi vào cơ sở dữ liệu mà nó sở hữu, mọi trao đổi dữ liệu
% chéo buộc phải đi qua API hoặc sự kiện được công bố công khai. Nguyên tắc này loại bỏ
% hoàn toàn sự kết dính schema giữa các service --- một trong những nguyên nhân hàng đầu khiến
% monolith khó phát triển độc lập.

% \textbf{Polyglot Persistence} --- Database per Service mở ra khả năng mỗi service chọn loại
% cơ sở dữ liệu phù hợp nhất với kiểu truy cập của miền mình thay vì ép tất cả vào một mô
% hình duy nhất. Trong hệ thống, lựa chọn được thực hiện dựa trên đặc thù truy cập:
% \textbf{PostgreSQL 17} (10 database riêng biệt) cho các service cần tính nhất quán mạnh,
% ràng buộc khóa ngoại và giao dịch ACID; \textbf{Redis 8} cho cart-service (key-value với
% TTL) và flash-sale-service (bộ đếm nguyên tử); \textbf{Elasticsearch 8.18} là kho đọc của
% search-service, tối ưu cho tìm kiếm toàn văn đa tiêu chí; \textbf{Apache Kafka KRaft} đóng
% vai trò message broker bền vững, đảm bảo giao tiếp bất đồng bộ tin cậy giữa các service.

% Database per Service tạo ra một hệ quả không thể tránh: khi một nghiệp vụ yêu cầu cập
% nhật dữ liệu trên nhiều service (ví dụ: đặt hàng đồng thời phải cập nhật order\_db,
% inventory\_db, voucher\_db và payment\_db), không thể dùng giao dịch ACID truyền thống
% để bao phủ tất cả vì mỗi database thuộc về một service độc lập. Đây là bài toán nhất
% quán phân tán~\cite{kleppmann2017}, được giải quyết bằng các mẫu thiết kế ở các phần tiếp theo.

\section{Phân rã hệ thống và Database per Service}
\label{section:2.2}

Hệ thống thương mại điện tử phức tạp không thể được mô hình hóa bằng một mô hình dữ liệu
thống nhất duy nhất, vì các miền nghiệp vụ khác nhau nhìn nhận cùng một khái niệm theo những
cách không tương thích. Ví dụ, khái niệm \textit{Đơn hàng} trong miền đặt hàng mang ý nghĩa
về thông tin khách hàng, danh sách sản phẩm và phương thức thanh toán; trong khi ở miền kho
hàng, chính đơn hàng đó lại là tập hợp các mặt hàng cần trừ tồn kho. Domain-Driven Design
(DDD) giải quyết mâu thuẫn này thông qua khái niệm \textbf{Bounded Context}~\cite{evans2003}:
một ranh giới tường minh mà bên trong đó, một mô hình nghiệp vụ cụ thể có ý nghĩa nhất quán
và không cần giải thích thêm. Áp dụng vào đề tài, mỗi Bounded Context được ánh xạ thành một
microservice sở hữu mô hình nghiệp vụ riêng và không chia sẻ lớp domain với service khác.

Ranh giới giữa các context được hiện thực hóa qua nguyên tắc \textit{Database per
Service}~\cite{richardson2018}: mỗi service chỉ được phép đọc và ghi vào cơ sở dữ liệu mà
nó sở hữu, mọi trao đổi dữ liệu chéo buộc phải đi qua API hoặc sự kiện được công bố công
khai. Nguyên tắc này loại bỏ hoàn toàn sự kết dính schema giữa các service --- một trong
những nguyên nhân hàng đầu khiến monolith khó phát triển độc lập.

\textbf{Polyglot Persistence.} Database per Service mở ra khả năng mỗi service lựa chọn loại
cơ sở dữ liệu phù hợp nhất với kiểu truy cập của miền mình, thay vì ép tất cả vào một mô
hình duy nhất. Tùy theo đặc thù truy cập, một service có thể dùng cơ sở dữ liệu quan hệ khi
cần tính nhất quán mạnh và giao dịch ACID, dùng kho key-value trong bộ nhớ cho các thao tác
đòi hỏi độ trễ thấp, hoặc dùng search engine cho nhu cầu tìm kiếm toàn văn. Việc lựa chọn và
lý giải cụ thể từng công nghệ trong hệ thống của đề tài được trình bày trong Chương~3.

Database per Service đồng thời tạo ra một hệ quả không thể tránh: khi một nghiệp vụ yêu cầu
cập nhật dữ liệu trên nhiều service --- chẳng hạn thao tác đặt hàng đồng thời tác động đến
dữ liệu đơn hàng, tồn kho, khuyến mãi và thanh toán --- không thể dùng một giao dịch ACID
truyền thống để bao phủ tất cả, vì mỗi cơ sở dữ liệu thuộc về một service độc lập. Đây chính
là bài toán nhất quán dữ liệu phân tán~\cite{kleppmann2017}, được giải quyết bằng các mẫu
thiết kế trình bày ở các phần tiếp theo của chương.
% =======================================================================
% \section{Saga Orchestration --- giao dịch phân tán nhiều bước}
% \label{section:2.3}

% Giao dịch ACID của cơ sở dữ liệu quan hệ đảm bảo rằng một tập hợp thao tác trên cùng một
% database hoặc cùng thành công hoặc cùng bị hủy. Tuy nhiên, giao thức Two-Phase Commit (2PC)
% --- cách mở rộng ACID sang nhiều database --- có hai vấn đề nghiêm trọng trong kiến trúc
% microservices: 2PC yêu cầu coordinator trung tâm giữ khóa trên tất cả participant trong suốt
% giao dịch, tạo điểm nghẽn hiệu năng và điểm lỗi đơn; đồng thời khả năng hỗ trợ 2PC phụ
% thuộc vào driver cơ sở dữ liệu và với Redis, Elasticsearch thì không tồn tại.

% \textbf{Saga Pattern}~\cite{garcia1987, richardson2018} là giải pháp thay thế: phân rã giao dịch phân tán lớn thành chuỗi các
% giao dịch cục bộ $T_1, T_2, \ldots, T_n$, mỗi giao dịch thực thi trong phạm vi một service
% duy nhất và commit ngay khi hoàn thành. Với mỗi $T_i$ thành công, hệ thống phải xác định
% giao dịch bù trừ $C_i$ có khả năng hoàn tác tác động của $T_i$; nếu $T_k$ thất bại, hệ
% thống thực thi $C_{k-1}, C_{k-2}, \ldots, C_1$ theo thứ tự ngược. Điểm khác biệt so với ACID
% là Saga không cung cấp \textit{isolation}: dữ liệu tạm thời không nhất quán giữa các service
% nhưng hội tụ về đúng sau khi chuỗi Saga hoàn thành (\textit{eventual consistency}).

% \subsection{Saga Choreography}

% Trong biến thể Choreography, không có thành phần điều phối trung tâm: mỗi service tự lắng
% nghe sự kiện từ Kafka và tự quyết định hành động kế tiếp dựa trên loại sự kiện nhận được.
% Ví dụ, nếu áp dụng Choreography cho luồng đặt hàng, inventory-service sẽ lắng nghe
% \texttt{order-created}, xử lý đặt chỗ tồn kho rồi phát \texttt{inventory-updated}; tiếp đó
% payment-service lắng nghe \texttt{inventory-updated} và tự tiến hành thanh toán mà không cần
% ai ra lệnh. Mỗi service chỉ biết sự kiện nó quan tâm, không biết gì về các service khác
% trong luồng --- tạo ra tính kết dính rất lỏng (loose coupling).

% Ưu điểm của Choreography là phù hợp với luồng đơn giản, ít bước, ít thay đổi: không cần
% thành phần trung tâm, mỗi service hoạt động hoàn toàn độc lập. Nhược điểm rõ ràng khi
% luồng phức tạp: trạng thái Saga bị phân tán khắp các service, không có điểm nào có bức
% tranh toàn cảnh về luồng đang ở bước nào; xử lý bù trừ trở nên khó kiểm soát vì không
% có orchestrator chủ động phát lệnh theo thứ tự ngược.

% \begin{figure}[H]
%     \centering
%     \includegraphics[width=0.85\textwidth]{Hinhve/chuong2/saga-choreography.png}
%     \caption{Saga Choreography Pattern --- mỗi service tự phản ứng theo sự kiện}
%     \label{fig:saga-choreography}
% \end{figure}

% \subsection{Saga Orchestration}

% Trong biến thể Orchestration, một thành phần trung tâm (orchestrator) nắm toàn bộ luồng và
% ra lệnh tuần tự đến từng service tham gia. Orchestrator biết chính xác Saga đang ở bước
% nào, quyết định bước tiếp theo và khi có sự cố, chủ động phát lệnh bù trừ theo thứ tự
% ngược.

% Ưu điểm của Orchestration là toàn bộ trạng thái Saga được tập trung tại một nơi nên dễ
% theo dõi, kiểm thử và mở rộng: thêm hoặc sửa một bước nghiệp vụ chỉ cần thay đổi
% orchestrator thay vì sửa nhiều service. Nhược điểm là orchestrator là thành phần phức tạp
% hơn và cần được thiết kế cẩn thận để không trở thành điểm nghẽn hay điểm lỗi đơn của toàn
% luồng.

% \begin{figure}[H]
%     \centering
%     \includegraphics[width=0.85\textwidth]{Hinhve/chuong2/saga-orchestrator.png}
%     \caption{Saga Orchestration Pattern --- orchestrator điều phối tập trung}
%     \label{fig:saga-orchestration}
% \end{figure}

% Sau khi đã phân tích riêng từng biến thể, Bảng~\ref{table:saga-comparison} đối chiếu trực
% tiếp Choreography và Orchestration theo các tiêu chí then chốt để làm cơ sở cho lựa chọn
% của hệ thống.

% \begin{table}[H]
% \centering
% \caption{So sánh Saga Choreography và Saga Orchestration}
% \label{table:saga-comparison}
% \begin{compacttable}
% \begin{tabularx}{\textwidth}{|L{2.8cm}|Y|Y|}
% \hline
% \textbf{Đặc điểm} & \textbf{Choreography} & \textbf{Orchestration} \\ \hline
% Cơ chế điều phối  & Mỗi service tự lắng nghe sự kiện và tự quyết định hành động kế tiếp & Orchestrator trung tâm ra lệnh tuần tự đến từng service \\ \hline
% Tính kết dính     & Lỏng --- các service không biết về nhau & Chặt hơn --- orchestrator nắm toàn bộ luồng \\ \hline
% Trạng thái Saga   & Phân tán khắp các service, khó theo dõi & Tập trung tại orchestrator, dễ kiểm tra \\ \hline
% Xử lý bù trừ     & Phức tạp --- khó xác định điểm kích hoạt bù trừ & Orchestrator chủ động phát lệnh bù trừ theo thứ tự ngược \\ \hline
% Thêm bước nghiệp vụ & Cần sửa nhiều service & Chỉ sửa orchestrator \\ \hline
% Phù hợp khi       & Luồng đơn, ít bước, ít thay đổi & Luồng phức tạp, nhiều bước, cần truy vết tập trung \\ \hline
% \end{tabularx}
% \end{compacttable}
% \end{table}

% Luồng đặt hàng trong hệ thống trải qua 6--8 bước liên quan đến năm service khác nhau
% (product, inventory, voucher, payment và notification) ngoài chính order-service điều phối,
% với hai nhánh rẽ phụ thuộc vào phương thức thanh toán (COD hoặc VNPAY) và nhiều điểm có thể
% thất bại. Orchestration là lựa chọn phù hợp hơn vì trạng thái Saga được tập trung tại
% order-service, logic bù trừ được xử lý một chỗ và thêm hoặc sửa bước nghiệp vụ chỉ cần
% thay đổi một nơi. Order-service đóng vai trò orchestrator duy nhất và duy trì vòng đời đơn
% hàng qua bốn trạng thái: \texttt{PENDING} (đang chờ đặt chỗ tồn kho), \texttt{STOCK\_RESERVED}
% (tồn kho đã đặt, chờ thanh toán), \texttt{CONFIRMED} (kết thúc thành công) và
% \texttt{CANCELLED} (thất bại, orchestrator phát lệnh bù trừ giải phóng tồn kho và hoàn trả
% voucher). Mỗi chuyển trạng thái được kích hoạt bởi sự kiện Kafka cụ thể, đảm bảo mọi đường
% chuyển đổi bất hợp lệ bị từ chối ở tầng code.

\section{Saga Orchestration --- giao dịch phân tán nhiều bước}
\label{section:2.3}

Giao dịch ACID của cơ sở dữ liệu quan hệ đảm bảo rằng một tập hợp thao tác trên cùng một
database hoặc cùng thành công hoặc cùng bị hủy. Tuy nhiên, giao thức Two-Phase Commit (2PC)
--- cách mở rộng ACID sang nhiều database --- bộc lộ hai vấn đề nghiêm trọng trong kiến trúc
microservices: 2PC yêu cầu coordinator trung tâm giữ khóa trên tất cả participant trong suốt
giao dịch, tạo điểm nghẽn hiệu năng và điểm lỗi đơn; đồng thời khả năng hỗ trợ 2PC phụ thuộc
vào driver cơ sở dữ liệu và không tồn tại với các kho dữ liệu như Redis hay Elasticsearch.

\textbf{Saga Pattern}~\cite{garcia1987, richardson2018} là giải pháp thay thế: phân rã giao
dịch phân tán lớn thành chuỗi các giao dịch cục bộ $T_1, T_2, \ldots, T_n$, mỗi giao dịch
thực thi trong phạm vi một service duy nhất và commit ngay khi hoàn thành. Với mỗi $T_i$
thành công, hệ thống phải xác định giao dịch bù trừ $C_i$ có khả năng hoàn tác tác động của
$T_i$; nếu $T_k$ thất bại, hệ thống thực thi $C_{k-1}, C_{k-2}, \ldots, C_1$ theo thứ tự
ngược. Điểm khác biệt so với ACID là Saga không cung cấp \textit{isolation}: dữ liệu có thể
tạm thời không nhất quán giữa các service, nhưng sẽ hội tụ về trạng thái đúng sau khi chuỗi
Saga hoàn thành (\textit{eventual consistency}).

\subsection{Saga Choreography}

Trong biến thể Choreography, không có thành phần điều phối trung tâm: mỗi service tự lắng
nghe sự kiện từ message broker và tự quyết định hành động kế tiếp dựa trên loại sự kiện nhận
được. Ví dụ, nếu áp dụng Choreography cho luồng đặt hàng, inventory-service sẽ lắng nghe
sự kiện \texttt{order-created}, xử lý đặt chỗ tồn kho rồi phát \texttt{inventory-updated};
tiếp đó payment-service lắng nghe \texttt{inventory-updated} và tự tiến hành thanh toán mà
không cần ai ra lệnh. Mỗi service chỉ biết những sự kiện nó quan tâm, không biết gì về các
service khác trong luồng --- tạo ra tính kết dính rất lỏng (loose coupling).

Ưu điểm của Choreography là phù hợp với luồng đơn giản, ít bước, ít thay đổi: không cần
thành phần trung tâm, mỗi service hoạt động hoàn toàn độc lập. Nhược điểm bộc lộ rõ khi
luồng trở nên phức tạp: trạng thái Saga bị phân tán khắp các service, không có điểm nào nắm
bức tranh toàn cảnh về luồng đang ở bước nào; việc xử lý bù trừ cũng trở nên khó kiểm soát
vì không có orchestrator chủ động phát lệnh theo thứ tự ngược.

\begin{figure}[H]
    \centering
    \includegraphics[width=0.85\textwidth]{Hinhve/chuong2/saga-choreography.png}
    \caption{Saga Choreography Pattern --- mỗi service tự phản ứng theo sự kiện}
    \label{fig:saga-choreography}
\end{figure}

\subsection{Saga Orchestration}

Trong biến thể Orchestration, một thành phần trung tâm (orchestrator) nắm toàn bộ luồng và
ra lệnh tuần tự đến từng service tham gia. Orchestrator biết chính xác Saga đang ở bước nào,
quyết định bước tiếp theo, và khi có sự cố sẽ chủ động phát lệnh bù trừ theo thứ tự ngược.

Ưu điểm của Orchestration là toàn bộ trạng thái Saga được tập trung tại một nơi nên dễ theo
dõi, kiểm thử và mở rộng: thêm hoặc sửa một bước nghiệp vụ chỉ cần thay đổi orchestrator
thay vì sửa nhiều service. Nhược điểm là orchestrator là thành phần phức tạp hơn và cần được
thiết kế cẩn thận để không trở thành điểm nghẽn hay điểm lỗi đơn của toàn luồng.

\begin{figure}[H]
    \centering
    \includegraphics[width=0.85\textwidth]{Hinhve/chuong2/saga-orchestrator.png}
    \caption{Saga Orchestration Pattern --- orchestrator điều phối tập trung}
    \label{fig:saga-orchestration}
\end{figure}

Sau khi đã phân tích riêng từng biến thể, Bảng~\ref{table:saga-comparison} đối chiếu trực
tiếp Choreography và Orchestration theo các tiêu chí then chốt để làm cơ sở cho lựa chọn
của hệ thống.

\begin{table}[H]
\centering
\caption{So sánh Saga Choreography và Saga Orchestration}
\label{table:saga-comparison}
\begin{compacttable}
\begin{tabularx}{\textwidth}{|L{2.8cm}|Y|Y|}
\hline
\textbf{Đặc điểm} & \textbf{Choreography} & \textbf{Orchestration} \\ \hline
Cơ chế điều phối  & Mỗi service tự lắng nghe sự kiện và tự quyết định hành động kế tiếp & Orchestrator trung tâm ra lệnh tuần tự đến từng service \\ \hline
Tính kết dính     & Lỏng --- các service không biết về nhau & Chặt hơn --- orchestrator nắm toàn bộ luồng \\ \hline
Trạng thái Saga   & Phân tán khắp các service, khó theo dõi & Tập trung tại orchestrator, dễ kiểm tra \\ \hline
Xử lý bù trừ     & Phức tạp --- khó xác định điểm kích hoạt bù trừ & Orchestrator chủ động phát lệnh bù trừ theo thứ tự ngược \\ \hline
Thêm bước nghiệp vụ & Cần sửa nhiều service & Chỉ sửa orchestrator \\ \hline
Phù hợp khi       & Luồng đơn, ít bước, ít thay đổi & Luồng phức tạp, nhiều bước, cần truy vết tập trung \\ \hline
\end{tabularx}
\end{compacttable}
\end{table}

Đối chiếu với đặc thù của đề tài, luồng đặt hàng là một quy trình nhiều bước, trải dài trên
nhiều service và có nhiều điểm có thể thất bại. Với mức độ phức tạp đó, Orchestration là lựa
chọn phù hợp hơn: trạng thái Saga được tập trung tại một orchestrator duy nhất, logic bù trừ
được xử lý ở một nơi, và việc thêm hoặc sửa bước nghiệp vụ chỉ tác động đến một thành phần.
Vì vậy, đề tài lựa chọn Saga Orchestration cho luồng đặt hàng; thiết kế chi tiết của
orchestrator, máy trạng thái đơn hàng và các luồng bù trừ tương ứng sẽ được trình bày trong
Chương~4.
% =======================================================================
% \section{Transactional Outbox và Idempotent Consumer}
% \label{section:2.4}

% Sau khi order-service tạo đơn hàng trong cơ sở dữ liệu, nó cần phát sự kiện
% \texttt{order-created} lên Kafka để inventory-service tiến hành đặt chỗ tồn kho. Cách
% tiếp cận thông thường là ghi dữ liệu vào database rồi ngay sau đó gọi \texttt{kafkaTemplate.send()}
% --- đây là vấn đề \textbf{dual-write}. Nếu service gặp sự cố sau khi commit transaction
% database nhưng trước khi kịp gọi Kafka, sự kiện sẽ không bao giờ được phát: đơn hàng tồn
% tại trong database với trạng thái \texttt{PENDING} nhưng inventory-service không nhận được
% yêu cầu đặt chỗ. Ngược lại, nếu Kafka nhận sự kiện nhưng database bị rollback, inventory
% sẽ đặt chỗ cho đơn hàng không tồn tại. Cả hai tình huống đều dẫn đến mất đồng bộ không
% phát hiện được bằng cách quan sát từng thành phần riêng lẻ.

% \textbf{Transactional Outbox Pattern}~\cite{richardson2018} giải quyết dual-write bằng cách khai thác ACID của
% database đã có: thay vì gọi Kafka trực tiếp, service ghi bản ghi sự kiện vào bảng
% \texttt{outbox} trong \textbf{cùng một giao dịch database} với dữ liệu nghiệp vụ. Một tiến
% trình nền (OutboxPoller) chạy định kỳ quét bảng outbox, đọc các sự kiện chưa gửi và phát
% lên Kafka; sau khi Kafka xác nhận, bản ghi được đánh dấu \texttt{processed = true}. Sơ đồ
% luồng được minh họa ở Hình~\ref{fig:outbox-pattern}. OutboxPoller sử dụng
% \texttt{SELECT ... FOR UPDATE SKIP LOCKED}: \texttt{FOR UPDATE} khóa các hàng được chọn
% ngăn instance khác lấy trùng; \texttt{SKIP LOCKED} bỏ qua hàng đã bị khóa thay vì chờ ---
% cho phép nhiều instance xử lý song song. Cơ chế này cung cấp \textbf{at-least-once delivery}:
% nếu service gặp sự cố sau khi Kafka nhận nhưng trước khi outbox được đánh dấu, Poller sẽ
% phát lại sự kiện sau khi khởi động lại.

% \begin{figure}[H]
%     \centering
%     \includegraphics[width=0.85\textwidth]{Hinhve/chuong2/outbox-pattern.png}
%     \caption{Luồng hoạt động của Transactional Outbox Pattern}
%     \label{fig:outbox-pattern}
% \end{figure}

% \textbf{Idempotent Consumer} đảm bảo rằng việc xử lý cùng một thông điệp nhiều lần không
% tạo ra tác động phụ khác biệt so với chỉ xử lý một lần. Cơ chế dựa trên bảng
% \texttt{processed\_events} với khóa chính là \texttt{event\_id}: consumer nhận thông điệp
% với \texttt{event\_id} duy nhất; trong cùng một transaction, cố gắng INSERT \texttt{event\_id}
% vào \texttt{processed\_events} và thực thi logic nghiệp vụ. Nếu INSERT thành công (event
% chưa xử lý), logic được commit; nếu INSERT vi phạm unique constraint (event đã xử lý),
% transaction rollback và thông điệp bị bỏ qua hoàn toàn. Kết hợp lại: Transactional Outbox
% đảm bảo sự kiện không bao giờ bị mất (at-least-once); Idempotent Consumer loại bỏ tác động
% trùng lặp. Sự kết hợp này tạo ra \textbf{hiệu ứng exactly-once processing} ở tầng ứng dụng,
% ngay cả khi hạ tầng Kafka chỉ đảm bảo at-least-once delivery.

\section{Transactional Outbox và Idempotent Consumer}
\label{section:2.4}

Khi một service vừa cập nhật dữ liệu nghiệp vụ vào database, vừa cần phát một sự kiện lên
message broker để các service khác phản ứng, nó phải thực hiện hai thao tác ghi trên hai hệ
thống khác nhau --- đây chính là vấn đề \textbf{dual-write}. Lấy ví dụ luồng đặt hàng: sau
khi tạo đơn trong database, order-service cần phát sự kiện \texttt{order-created} để
inventory-service đặt chỗ tồn kho. Nếu service gặp sự cố sau khi đã commit transaction
database nhưng trước khi kịp phát sự kiện, sự kiện sẽ không bao giờ được gửi đi: đơn hàng
tồn tại trong database nhưng inventory-service không hề nhận được yêu cầu đặt chỗ. Ngược lại,
nếu sự kiện đã được phát nhưng transaction database bị rollback, inventory-service sẽ đặt chỗ
cho một đơn hàng không tồn tại. Cả hai tình huống đều dẫn đến mất đồng bộ dữ liệu giữa các
service mà không thể phát hiện bằng cách quan sát từng thành phần riêng lẻ.

\textbf{Transactional Outbox Pattern}~\cite{richardson2018} giải quyết dual-write bằng cách
khai thác chính tính ACID sẵn có của database: thay vì phát sự kiện trực tiếp lên broker,
service ghi bản ghi sự kiện vào một bảng \texttt{outbox} trong \textbf{cùng một giao dịch
database} với dữ liệu nghiệp vụ. Nhờ đó, hoặc cả dữ liệu nghiệp vụ lẫn sự kiện cùng được lưu,
hoặc cả hai cùng bị hủy --- không thể xảy ra tình trạng một trong hai bị mất. Một tiến trình
nền (outbox poller) chạy định kỳ quét bảng outbox, đọc các sự kiện chưa gửi và phát lên
broker; sau khi broker xác nhận, bản ghi tương ứng được đánh dấu là đã xử lý. Sơ đồ luồng
được minh họa ở Hình~\ref{fig:outbox-pattern}.

\begin{figure}[H]
    \centering
    \includegraphics[width=0.85\textwidth]{Hinhve/chuong2/outbox-pattern.png}
    \caption{Luồng hoạt động của Transactional Outbox Pattern}
    \label{fig:outbox-pattern}
\end{figure}

Để nhiều instance của cùng một service có thể quét bảng outbox song song mà không xử lý
trùng, poller sử dụng truy vấn \texttt{SELECT ... FOR UPDATE SKIP LOCKED}: mệnh đề
\texttt{FOR UPDATE} khóa các hàng được chọn để ngăn instance khác lấy trùng, còn
\texttt{SKIP LOCKED} cho phép bỏ qua các hàng đang bị khóa thay vì phải chờ. Cơ chế này
mang lại bảo đảm \textbf{at-least-once delivery}: nếu service gặp sự cố sau khi broker đã
nhận sự kiện nhưng trước khi bản ghi outbox kịp được đánh dấu, poller sẽ phát lại sự kiện
đó sau khi khởi động lại.

\textbf{Idempotent Consumer Pattern} đảm bảo rằng việc xử lý cùng một thông điệp nhiều lần
không tạo ra tác động khác biệt so với chỉ xử lý một lần --- một yêu cầu tất yếu khi tầng
vận chuyển chỉ bảo đảm at-least-once và do đó có thể giao cùng một thông điệp nhiều lần. Cơ
chế thường dựa trên một bảng lưu lại định danh của những sự kiện đã xử lý (\texttt{event\_id}):
khi nhận một thông điệp, consumer kiểm tra xem định danh của nó đã tồn tại hay chưa; nếu chưa,
consumer thực thi logic nghiệp vụ và ghi nhận định danh đó trong cùng một giao dịch; nếu đã
tồn tại, thông điệp được xem là trùng lặp và bị bỏ qua hoàn toàn.

Kết hợp hai pattern: Transactional Outbox bảo đảm sự kiện không bao giờ bị mất (at-least-once),
còn Idempotent Consumer loại bỏ tác động của các sự kiện trùng lặp. Sự kết hợp này tạo ra
\textbf{hiệu ứng exactly-once processing} ở tầng ứng dụng, ngay cả khi hạ tầng message broker
bên dưới chỉ bảo đảm at-least-once delivery.
% =======================================================================
% \section{Xử lý đồng thời cao --- Redis Lua atomic và mô hình phòng thủ đa lớp}
% \label{section:2.5}

% Flash sale tạo ra tình huống đồng thời cực cao: hàng trăm đến hàng nghìn người dùng đồng
% thời cố gắng mua một số lượng hàng giới hạn. Vấn đề nằm ở \textbf{race condition kinh điển}:
% giả sử còn 1 slot, hai request A và B đồng thời đọc stock~=~1, cả hai thấy còn hàng, cả hai
% tiến hành mua --- kết quả: 1 slot được bán cho 2 người (over-sell). Hai cách giải quyết
% truyền thống trong cơ sở dữ liệu quan hệ đều không phù hợp cho tải cao: pessimistic lock
% (\texttt{FOR UPDATE}) tạo hàng đợi tuần tự, hiệu năng sụp đổ khi có 500 VU đồng thời;
% optimistic lock (version column) tỷ lệ retry tăng tương ứng số VU, làm tăng tải thay vì
% giảm.

% \textbf{Redis single-threaded model} cung cấp nguyên tử tự nhiên: mỗi lệnh được thực thi
% hoàn chỉnh trước khi lệnh tiếp theo bắt đầu, không có sự xen ngang giữa các lệnh từ bất kỳ
% kết nối nào. Tuy nhiên, bài toán flash sale đòi hỏi thực hiện \textit{nhiều} lệnh như một
% khối không thể tách rời: kiểm tra người dùng đã mua chưa, kiểm tra stock, giảm stock và ghi
% nhận người mua --- nếu thực thi tuần tự các lệnh riêng biệt, khoảng trống giữa các lệnh vẫn
% tạo ra cửa sổ cho race condition.

% \textbf{Redis Lua Script}~\cite{redis_docs} giải quyết triệt để: khi client gửi script Lua lên Redis, server
% thực thi toàn bộ script như một khối nguyên tử duy nhất --- không một lệnh nào khác được
% thực thi trong suốt thời gian script chạy. Script flash sale thực hiện đầy đủ trong một lần
% gọi: (1)~kiểm tra \texttt{flash-sale:\{id\}:buyers} xem \texttt{user\_id} đã mua chưa ---
% nếu rồi trả về \texttt{-1}; (2)~đọc \texttt{flash-sale:\{id\}:stock} --- nếu $\le 0$ trả
% về \texttt{0}; (3)~\texttt{DECR} stock và \texttt{SADD} user\_id vào buyers set, trả về
% \texttt{1}. Không có bất kỳ khoảng trống nào giữa các bước, không có race condition nào
% có thể xảy ra kể cả khi có hàng nghìn request đồng thời.

% \textbf{Mô hình phòng thủ đa lớp} bổ sung cho Lua script để bảo vệ trước sự cố hạ tầng:
% \textit{Lớp 1 --- Nguyên tử Redis}: Lua script xử lý phần lớn request trong bộ nhớ với độ
% trễ dưới mili-giây, loại bỏ ngay các request không đủ điều kiện trước khi đến tầng database
% hay Kafka. \textit{Lớp 2 --- Kafka Compensation}: nếu publish Kafka thất bại sau khi Lua
% script xác nhận slot, service lập tức thực thi compensation: \texttt{INCR} stock và
% \texttt{SREM} user\_id khỏi buyers set. \textit{Lớp 3 --- ReconciliationScheduler}: chạy
% mỗi 5 phút, đối chiếu \texttt{sold\_count} trong database với số đơn flash-sale thực tế;
% nếu phát hiện drift, bù trừ counter Redis về giá trị đúng. Ba lớp hoạt động độc lập và bổ
% sung cho nhau: lớp 1 nhanh nhưng không bảo vệ khỏi lỗi hạ tầng; lớp 2 xử lý lỗi ngay lập
% tức; lớp 3 bắt sai lệch tích lũy mà hai lớp trên bỏ sót.

\section{Xử lý đồng thời cao --- Redis Lua atomic và mô hình phòng thủ đa lớp}
\label{section:2.5}

Flash sale tạo ra tình huống đồng thời cực cao: hàng trăm đến hàng nghìn người dùng cùng lúc
cố gắng mua một số lượng hàng giới hạn. Vấn đề cốt lõi nằm ở \textbf{race condition kinh
điển}: giả sử chỉ còn một slot, hai request A và B đồng thời đọc thấy stock~=~1, cả hai đều
kết luận là còn hàng và cùng tiến hành mua --- kết quả là một slot bị bán cho hai người
(over-sell). Hai cách giải quyết truyền thống trong cơ sở dữ liệu quan hệ đều không phù hợp
với điều kiện tải cao: khóa bi quan (pessimistic lock) biến các request thành một hàng đợi
tuần tự khiến hiệu năng sụp đổ khi số lượng truy cập đồng thời lớn; còn khóa lạc quan
(optimistic lock) lại khiến tỷ lệ retry tăng theo mức độ đồng thời, làm tăng tải thay vì
giảm.

\textbf{Mô hình đơn luồng của Redis.} Redis xử lý các lệnh theo mô hình single-threaded: mỗi
lệnh được thực thi trọn vẹn trước khi lệnh tiếp theo bắt đầu, không có sự xen ngang giữa các
lệnh đến từ bất kỳ kết nối nào. Đặc tính này cung cấp tính nguyên tử tự nhiên ở mức từng
lệnh. Tuy nhiên, bài toán flash sale đòi hỏi thực hiện \textit{nhiều} lệnh như một khối không
thể tách rời --- kiểm tra người dùng đã mua hay chưa, kiểm tra tồn kho, giảm tồn kho và ghi
nhận người mua. Nếu thực thi các lệnh này một cách riêng lẻ, khoảng trống giữa chúng vẫn tạo
ra cửa sổ cho race condition xảy ra.

\textbf{Redis Lua Script}~\cite{redis_docs} giải quyết triệt để vấn đề trên: khi client gửi
một script Lua lên Redis, server thực thi toàn bộ script như một khối nguyên tử duy nhất ---
không lệnh nào khác được thực thi xen vào trong suốt thời gian script chạy. Nhờ đó, toàn bộ
chuỗi thao tác kiểm tra điều kiện và trừ tồn kho được gộp vào một lần gọi và diễn ra liền
mạch, không để lại bất kỳ khoảng trống nào cho race condition, kể cả khi có hàng nghìn
request ập đến đồng thời. Đây là nền tảng cho lớp phòng thủ đầu tiên chống over-sell.

\textbf{Mô hình phòng thủ đa lớp.} Bản thân Lua script chỉ bảo đảm tính nguyên tử trong bộ
nhớ Redis; nó không tự bảo vệ hệ thống trước các sự cố ở những thành phần hạ tầng khác như
database hay message broker. Vì vậy, giải pháp chống over-sell hoàn chỉnh được tổ chức theo
mô hình phòng thủ ba lớp bổ sung cho nhau:

\begin{itemize}[noitemsep]
    \item \textbf{Lớp 1 --- Nguyên tử trong bộ nhớ:} Lua script xử lý phần lớn request ngay
    trong bộ nhớ Redis với độ trễ rất thấp, loại bỏ sớm các request không đủ điều kiện trước
    khi chúng chạm tới tầng database hay message broker.
    \item \textbf{Lớp 2 --- Bù trừ tức thời:} khi một thao tác ở tầng hạ tầng thất bại sau
    khi slot đã được xác nhận ở Lớp 1, hệ thống lập tức thực hiện thao tác bù trừ để hoàn lại
    slot, tránh tình trạng slot bị giữ nhưng đơn hàng không được tạo.
    \item \textbf{Lớp 3 --- Đối chiếu định kỳ:} một tác vụ nền chạy theo chu kỳ đối chiếu số
    liệu tồn kho trên Redis với số đơn flash sale thực tế trong database, phát hiện và sửa
    các sai lệch tích lũy mà hai lớp trên có thể bỏ sót.
\end{itemize}

Ba lớp hoạt động độc lập nhưng bổ sung cho nhau: Lớp 1 nhanh nhưng không chống được lỗi hạ
tầng; Lớp 2 xử lý lỗi ngay tại thời điểm phát sinh; Lớp 3 bắt các sai lệch còn sót lại theo
thời gian. Thiết kế chi tiết của từng lớp, cùng cấu trúc khóa Redis và các tác vụ nền tương
ứng, sẽ được trình bày trong Chương~4.
% =======================================================================
% \section{Bảo mật tập trung}
% \label{section:2.6}

% JSON Web Token (JWT)~\cite{rfc7519} là tiêu chuẩn mở cho phép truyền thông tin định danh đã được ký số
% giữa các bên mà không cần tra cứu trạng thái session trên server. Hệ thống sử dụng
% \textbf{RS256} (RSA + SHA-256) thay vì HS256 (HMAC): với HS256, cả bên ký lẫn bên xác minh
% phải chia sẻ cùng một khóa bí mật --- không phù hợp với kiến trúc nhiều service vì mỗi
% service phải biết secret key, tạo rủi ro bảo mật nếu một service bị xâm phạm. Với RS256,
% Keycloak giữ private key để ký token; các service chỉ cần public key (lấy từ JWKS endpoint)
% để xác minh chữ ký --- private key không bao giờ rời khỏi Keycloak.

% Trong kiến trúc microservices, nếu mỗi service tự xác minh JWT với Keycloak, mỗi request
% sinh thêm một lời gọi mạng từ từng service, tạo bottleneck và điểm lỗi đơn. \textbf{Trusted
% Subsystem Pattern} giải quyết bằng cách chỉ định một điểm thực thi xác thực duy nhất ---
% API Gateway --- và cho phép tất cả service nội bộ tin tưởng thông tin định danh do Gateway
% cung cấp mà không xác minh lại. Luồng cụ thể:

% \begin{enumerate}[noitemsep]
%     \item Client gửi request kèm Access Token trong header \texttt{Authorization: Bearer}.
%     \item API Gateway \textbf{loại bỏ} (strip) hoàn toàn các header \texttt{X-User-Id},
%           \texttt{X-User-Roles}, \texttt{X-User-Email} từ request gốc của client --- ngăn
%           chặn mọi nỗ lực giả mạo định danh từ phía client.
%     \item Gateway xác minh chữ ký JWT bằng public key lấy từ JWKS endpoint của Keycloak,
%           kiểm tra \texttt{exp} (hạn token) và \texttt{iss} (nhà phát hành).
%     \item Sau khi xác minh thành công, Gateway \textbf{chèn} thông tin định danh đã được
%           xác thực vào ba header: \texttt{X-User-Id} (sub claim), \texttt{X-User-Roles}
%           (realm\_access.roles), \texttt{X-User-Email} (email claim).
%     \item Request được chuyển tiếp đến service đích; service đọc trực tiếp các header này
%           mà không cần JWT library.
% \end{enumerate}

% \begin{figure}[H]
%     \centering
%     \includegraphics[width=0.85\textwidth]{Hinhve/chuong4/keycloak.png}
%     \caption{Luồng xác thực Trusted Subsystem Pattern: Gateway xác minh JWT và inject header định danh}
%     \label{fig:trusted-subsystem}
% \end{figure}

% Pattern này chỉ an toàn khi kết hợp với cách ly mạng tầng hạ tầng: nếu service backend
% có thể bị gọi trực tiếp từ bên ngoài, kẻ tấn công có thể bỏ qua Gateway và tự chèn header
% tùy ý. Trong hệ thống, tất cả service business (cổng 8081--8093) chỉ kết nối vào mạng
% \texttt{internal} của Docker, không expose cổng ra host; API Gateway (cổng 8080) là điểm
% duy nhất kết nối cả mạng \texttt{internal} lẫn \texttt{public}, không có đường nào để client
% gọi trực tiếp vào business service mà không qua Gateway.

\section{Bảo mật tập trung}
\label{section:2.6}

JSON Web Token (JWT)~\cite{rfc7519} là tiêu chuẩn mở cho phép truyền thông tin định danh đã
được ký số giữa các bên mà không cần tra cứu trạng thái session trên server. Một lựa chọn
quan trọng khi sử dụng JWT là thuật toán ký. Với HS256 (HMAC), bên ký và bên xác minh phải
chia sẻ cùng một khóa bí mật --- điều này không phù hợp với kiến trúc nhiều service, bởi mỗi
service đều phải nắm giữ khóa bí mật, kéo theo rủi ro lan rộng nếu chỉ một service bị xâm
phạm. Ngược lại, với \textbf{RS256} (RSA kết hợp SHA-256), nhà cung cấp định danh giữ private
key để ký token, còn các service chỉ cần public key để xác minh chữ ký; private key không bao
giờ rời khỏi nơi phát hành. Vì vậy RS256 là lựa chọn phù hợp cho môi trường phân tán.

Trong kiến trúc microservices, nếu mỗi service tự xác minh JWT một cách độc lập với nhà cung
cấp định danh thì mỗi request có thể phát sinh thêm các lời gọi mạng, gây tốn kém và tạo điểm
nghẽn. \textbf{Trusted Subsystem Pattern} giải quyết vấn đề này bằng cách chỉ định một điểm
thực thi xác thực duy nhất --- thường là API Gateway --- và cho phép các service nội bộ tin
tưởng thông tin định danh do điểm này cung cấp mà không cần xác minh lại. Luồng xử lý tổng
quát như sau:

\begin{enumerate}[noitemsep]
    \item Client gửi request kèm Access Token trong header \texttt{Authorization: Bearer}.
    \item API Gateway loại bỏ (strip) hoàn toàn các header định danh có trong request gốc của
          client, nhằm ngăn chặn mọi nỗ lực giả mạo định danh từ phía client.
    \item Gateway xác minh chữ ký JWT bằng public key của nhà cung cấp định danh, đồng thời
          kiểm tra hạn sử dụng và nhà phát hành của token.
    \item Sau khi xác minh thành công, Gateway chèn các thông tin định danh đã được xác thực
          (định danh người dùng, vai trò, email) vào những header tin cậy.
    \item Request được chuyển tiếp đến service đích; service đọc trực tiếp các header này mà
          không cần tự xử lý JWT.
\end{enumerate}

\begin{figure}[H]
    \centering
    \includegraphics[width=0.85\textwidth]{Hinhve/chuong4/keycloak.png}
    \caption{Luồng xác thực Trusted Subsystem Pattern: Gateway xác minh JWT và inject header định danh}
    \label{fig:trusted-subsystem}
\end{figure}

Cần lưu ý rằng pattern này chỉ thực sự an toàn khi đi kèm với cơ chế cách ly mạng ở tầng hạ
tầng. Nếu các service backend có thể bị gọi trực tiếp từ bên ngoài, kẻ tấn công hoàn toàn có
thể bỏ qua Gateway và tự chèn các header định danh tùy ý để mạo danh người dùng. Do đó, điều
kiện tiên quyết để Trusted Subsystem hoạt động an toàn là các service nội bộ phải được đặt
trong một mạng riêng, không phơi bày ra bên ngoài, sao cho mọi request từ client bắt buộc
phải đi qua Gateway. Cách hệ thống của đề tài hiện thực hóa sự cách ly mạng này được trình
bày trong Chương~5.
% =======================================================================
% \section{Các mẫu thiết kế hỗ trợ: CQRS-lite, Circuit Breaker và Observability}
% \label{section:2.7}

% \subsection{CQRS-lite --- tách biệt mô hình đọc và ghi cho tìm kiếm}

% CQRS (Command Query Responsibility Segregation)~\cite{fowler2011cqrs} là mẫu thiết kế tách hoàn toàn mô hình
% ghi (Command) và mô hình đọc (Query) của hệ thống, mỗi mô hình được lưu trữ trên kho dữ
% liệu tối ưu cho kiểu thao tác của mình. Hệ thống áp dụng biến thể \textbf{CQRS-lite}
% cho chức năng tìm kiếm sản phẩm:

% \begin{itemize}[noitemsep]
%     \item \textbf{Mô hình ghi}: product-service quản lý catalog trên PostgreSQL với schema
%           quan hệ chuẩn hóa, đảm bảo tính nhất quán và toàn vẹn dữ liệu (khóa ngoại,
%           constraints, giao dịch ACID). Đây là nguồn dữ liệu chân thực (source of truth).
%     \item \textbf{Mô hình đọc}: search-service duy trì một bản sao denormalized trong
%           Elasticsearch index \texttt{products}, được tối ưu cho full-text search, lọc đa
%           tiêu chí (category, brand, price range) và xếp hạng mức độ liên quan.
%     \item \textbf{Đồng bộ bất đồng bộ}: khi quản trị viên tạo/sửa/xóa sản phẩm,
%           product-service phát sự kiện \texttt{product-created/updated/deleted} lên Kafka.
%           Search-service tiêu thụ sự kiện và cập nhật Elasticsearch, chấp nhận độ trễ
%           đồng bộ ngắn (thường dưới 1 giây) để đổi lấy hiệu năng tìm kiếm vượt trội.
% \end{itemize}

% \begin{figure}[H]
%     \centering
%     \includegraphics[width=0.85\textwidth]{Hinhve/chuong2/CQRS.png}
%     \caption{Kiến trúc CQRS-lite: PostgreSQL làm mô hình ghi, Elasticsearch làm mô hình đọc}
%     \label{fig:CQRS}
% \end{figure}

% Lý do lựa chọn CQRS-lite thay vì query trực tiếp PostgreSQL: full-text search với phân
% tích ngữ nghĩa (stemming, tokenization), lọc kết hợp nhiều trường và xếp hạng mức độ liên
% quan là những khả năng mà Elasticsearch cung cấp tự nhiên nhưng sẽ cần query SQL cực kỳ
% phức tạp (LIKE, GIN index, tsvector) để mô phỏng trên PostgreSQL với hiệu năng kém hơn
% đáng kể.

% \subsection{Circuit Breaker --- ngăn sự cố dây chuyền giữa các service}

% Khi một service gọi đồng bộ đến service khác qua OpenFeign và service đích đang gặp sự
% cố hoặc phản hồi chậm, nếu không có cơ chế bảo vệ, thread pool của service gọi sẽ bị
% chiếm giữ bởi các request đang chờ, dẫn đến cạn kiệt tài nguyên và sự cố lan truyền sang
% các service lân cận --- hiện tượng \textit{cascading failure}.

% Circuit Breaker (Resilience4j)~\cite{nygard2018, resilience4j_docs} bảo vệ bằng cách giám sát tỷ lệ lỗi của mỗi Feign call
% và hoạt động theo ba trạng thái:

% \begin{figure}[H]
%     \centering
%     \includegraphics[width=0.85\textwidth]{Hinhve/chuong2/circuit-breaker.png}
%     \caption{Sơ đồ ba trạng thái Circuit Breaker: CLOSED, OPEN và HALF\_OPEN}
%     \label{fig:circuit-breaker}
% \end{figure}

% \begin{itemize}[noitemsep]
%     \item \textbf{CLOSED}: mọi request được chuyển tiếp bình thường. Nếu tỷ lệ lỗi vượt
%           ngưỡng (50\% trong 10 lần gọi gần nhất), circuit chuyển sang OPEN.
%     \item \textbf{OPEN}: tất cả request bị từ chối ngay lập tức (fail-fast) mà không gọi
%           thật đến service đích, trả về kết quả fallback. Sau 10 giây, circuit chuyển
%           sang HALF\_OPEN.
%     \item \textbf{HALF\_OPEN}: cho phép một số lượng hạn chế request thật đến service
%           đích để kiểm tra. Nếu thành công, circuit trở về CLOSED; nếu thất bại, quay
%           lại OPEN.
% \end{itemize}

% Trong hệ thống, Circuit Breaker được cấu hình cho toàn bộ bảy lời gọi Feign:
% cart-service → product-service (xác minh sản phẩm khi thêm giỏ), order-service →
% product-service (snapshot giá), order-service → voucher-service (reserve/commit/release),
% order-service → cart-service (xóa giỏ sau khi đặt hàng), payment-service → order-service
% (lấy payment context), review-service → order-service (xác minh lịch sử mua) và
% flash-sale-service → order-service (đếm đơn phục vụ reconciliation). Cần lưu ý các tương
% tác trong luồng Saga như đặt chỗ tồn kho không dùng Feign mà đi qua sự kiện Kafka, nên
% không nằm trong phạm vi bảo vệ của Circuit Breaker. Kết hợp với Retry (tối đa 3 lần, delay
% 500ms) và Time Limiter (timeout 3 giây), Resilience4j đảm bảo sự cố tại một service không
% làm tê liệt service phụ thuộc.

% \textbf{Rate Limiting} tại API Gateway giới hạn số lượng request mỗi client có thể gửi
% trong một khoảng thời gian, bảo vệ hệ thống khỏi lạm dụng và đảm bảo công bằng giữa người
% dùng. Thuật toán Token Bucket được sử dụng: mỗi client được gán "thùng" với sức chứa tối đa
% $B$ token, token bổ sung với tốc độ $r$ token/giây, mỗi request tiêu thụ 1 token ---
% cho phép burst ngắn hạn nhưng giới hạn tốc độ trung bình dài hạn. Hai endpoint được giới
% hạn đặc biệt: \texttt{POST /api/orders} (10 req/s, burst 20) và
% \texttt{POST /api/flash-sales/*/purchase} (3 req/s, burst 5). Trạng thái token bucket lưu
% trên Redis, cho phép rate limiting nhất quán khi Gateway chạy nhiều instance.

% \subsection{Khả năng quan sát --- ba trụ cột giám sát hệ thống phân tán}

% Khi một request thất bại trong hệ thống 16 service, không thể biết lỗi phát sinh từ
% đâu mà không có công cụ quan sát đúng. Observability được xây dựng trên ba trụ cột bổ
% sung lẫn nhau:

% \textbf{Metrics (Prometheus + Grafana)}: thu thập định kỳ các giá trị số phản ánh tình
% trạng hệ thống --- request rate, latency percentile (p50/p95/p99), error rate, JVM memory,
% Kafka consumer lag. Metrics hiệu quả cho cảnh báo ngưỡng và lập kế hoạch năng lực, nhưng
% không giải thích nguyên nhân của từng sự kiện cụ thể.

% \textbf{Distributed Tracing (Zipkin)}: gán một \texttt{TraceId} duy nhất cho mỗi request,
% truyền qua tất cả service mà request đi qua. Mỗi đoạn xử lý tại một service tạo ra một
% \texttt{Span} với thời gian bắt đầu, kết thúc và các nhãn mô tả. Tập hợp tất cả Span của
% một Trace tái tạo chính xác hành trình của request qua hệ thống, xác định service nào
% đang gây độ trễ. Micrometer Tracing tự động propagate context qua HTTP header
% (\texttt{traceparent} theo chuẩn W3C) và qua Kafka message header.

% \textbf{Log tương quan theo trace context}: toàn bộ service dùng chung một Logback pattern
% thống nhất, nhúng tên service cùng \texttt{traceId} và \texttt{spanId} vào mỗi dòng log
% (các giá trị này được Micrometer Tracing tự động điền vào MDC). Nhờ \texttt{traceId} xuất
% hiện đồng nhất trên mọi dòng log của cùng một request, có thể lọc và ghép toàn bộ log liên
% quan đến một request xuyên suốt các service, đồng thời đối chiếu trực tiếp giữa dòng log và
% trace tương ứng trên Zipkin khi điều tra sự cố.

% \textbf{Scheduled Jobs và Distributed Lock} --- nhiều tác vụ phải chạy định kỳ theo thời
% gian: OutboxPoller phát sự kiện mỗi 1 giây, ReservationExpiryScheduler hủy đơn VNPAY hết
% hạn mỗi 60 giây, ReconciliationScheduler đối chiếu số liệu flash sale mỗi 5 phút. Khi
% service chạy nhiều instance, cần đảm bảo chỉ một instance thực thi scheduler tại mỗi thời
% điểm. Giải pháp là Distributed Lock trên Redis với \texttt{SET NX EX}: \texttt{NX} đảm bảo
% chỉ một instance giành được khóa; \texttt{EX} đảm bảo khóa tự động giải phóng nếu instance
% đang giữ khóa gặp sự cố, tránh deadlock vĩnh viễn. Đối với OutboxPoller, \texttt{FOR UPDATE
% SKIP LOCKED} đóng vai trò distributed lock ở tầng database: mỗi instance xử lý một tập con
% bản ghi outbox mà không xung đột, tối đa hóa thông lượng khi scale.

% % \begin{figure}[!htbp]
% %     \centering
% %     \fbox{\parbox{0.85\textwidth}{\centering\vspace{3cm}%
% %         \texttt{Hinhve/chuong2/observability.png}%
% %     \vspace{3cm}}}
% %     \caption{Ba trụ cột Observability: Metrics, Distributed Tracing và Log tương quan theo trace context}
% %     \label{fig:observability}
% % \end{figure}

\section{Các mẫu thiết kế hỗ trợ: CQRS-lite, Circuit Breaker và Observability}
\label{section:2.7}

\subsection{CQRS-lite --- tách biệt mô hình đọc và ghi cho tìm kiếm}

CQRS (Command Query Responsibility Segregation)~\cite{fowler2011cqrs} là mẫu thiết kế tách
biệt mô hình ghi (Command) và mô hình đọc (Query) của hệ thống, mỗi mô hình được lưu trữ
trên kho dữ liệu tối ưu cho kiểu thao tác của mình. Đề tài áp dụng một biến thể đơn giản hóa
là \textbf{CQRS-lite} cho chức năng tìm kiếm sản phẩm, dựa trên ba thành phần:

\begin{itemize}[noitemsep]
    \item \textbf{Mô hình ghi}: một service quản lý catalog sản phẩm trên cơ sở dữ liệu quan
          hệ với schema chuẩn hóa, bảo đảm tính nhất quán và toàn vẹn dữ liệu. Đây là nguồn
          dữ liệu chân thực (source of truth).
    \item \textbf{Mô hình đọc}: một service tìm kiếm duy trì một bản sao đã được phi chuẩn
          hóa (denormalized) trên search engine, tối ưu cho tìm kiếm toàn văn, lọc đa tiêu
          chí và xếp hạng mức độ liên quan.
    \item \textbf{Đồng bộ bất đồng bộ}: mỗi khi dữ liệu sản phẩm thay đổi, mô hình ghi phát
          sự kiện tương ứng lên message broker; mô hình đọc tiêu thụ sự kiện và cập nhật lại
          chỉ mục tìm kiếm, chấp nhận một độ trễ đồng bộ rất ngắn để đổi lấy hiệu năng đọc
          vượt trội.
\end{itemize}

\begin{figure}[H]
    \centering
    \includegraphics[width=0.85\textwidth]{Hinhve/chuong2/CQRS.png}
    \caption{Kiến trúc CQRS-lite: cơ sở dữ liệu quan hệ làm mô hình ghi, search engine làm mô hình đọc}
    \label{fig:CQRS}
\end{figure}

Lý do lựa chọn CQRS-lite thay vì truy vấn trực tiếp trên cơ sở dữ liệu quan hệ là: tìm kiếm
toàn văn kèm phân tích ngữ nghĩa, lọc kết hợp nhiều trường và xếp hạng mức độ liên quan là
những khả năng mà một search engine cung cấp sẵn, nhưng sẽ rất phức tạp và kém hiệu quả nếu
phải mô phỏng bằng truy vấn trên cơ sở dữ liệu quan hệ. Việc tách mô hình đọc ra một kho dữ
liệu chuyên dụng giúp chức năng tìm kiếm đạt hiệu năng cao mà không ảnh hưởng đến mô hình ghi.

\subsection{Circuit Breaker --- ngăn sự cố dây chuyền giữa các service}

Khi một service gọi đồng bộ đến service khác và service đích đang gặp sự cố hoặc phản hồi
chậm, nếu không có cơ chế bảo vệ, các luồng (thread) của service gọi sẽ bị chiếm giữ bởi
những request đang chờ, dẫn tới cạn kiệt tài nguyên và lan truyền sự cố sang các service lân
cận --- hiện tượng \textit{cascading failure}.

Circuit Breaker là mẫu thiết kế ngăn chặn hiện tượng này~\cite{nygard2018, resilience4j_docs}
bằng cách giám sát tỷ lệ lỗi của mỗi lời gọi và hoạt động theo ba trạng thái:

\begin{figure}[H]
    \centering
    \includegraphics[width=0.85\textwidth]{Hinhve/chuong2/circuit-breaker.png}
    \caption{Sơ đồ ba trạng thái Circuit Breaker: CLOSED, OPEN và HALF\_OPEN}
    \label{fig:circuit-breaker}
\end{figure}

\begin{itemize}[noitemsep]
    \item \textbf{CLOSED}: mọi request được chuyển tiếp bình thường. Khi tỷ lệ lỗi vượt quá
          ngưỡng định trước, circuit chuyển sang trạng thái OPEN.
    \item \textbf{OPEN}: mọi request bị từ chối ngay lập tức (fail-fast) mà không thực sự
          gọi đến service đích, thay vào đó trả về kết quả dự phòng (fallback). Sau một
          khoảng thời gian chờ, circuit chuyển sang HALF\_OPEN.
    \item \textbf{HALF\_OPEN}: cho phép một số lượng hạn chế request thật đi qua để thăm dò.
          Nếu chúng thành công, circuit trở về CLOSED; nếu thất bại, circuit quay lại OPEN.
\end{itemize}

Circuit Breaker thường được kết hợp với một số cơ chế khả năng phục hồi (resilience) khác
như Retry (tự động thử lại sau một khoảng trễ) và Time Limiter (giới hạn thời gian chờ phản
hồi), nhằm đảm bảo sự cố tại một service không làm tê liệt các service phụ thuộc. Cần lưu ý
rằng Circuit Breaker chỉ áp dụng cho các lời gọi đồng bộ; những tương tác bất đồng bộ qua
message broker không nằm trong phạm vi bảo vệ của nó, mà dựa vào các cơ chế tin cậy riêng
như đã trình bày ở mục Transactional Outbox. Việc cấu hình Circuit Breaker cho từng lời gọi
cụ thể trong hệ thống được trình bày trong Chương~4.

\textbf{Rate Limiting} là cơ chế giới hạn số lượng request mà mỗi client có thể gửi trong
một khoảng thời gian, qua đó bảo vệ hệ thống khỏi bị lạm dụng và bảo đảm sự công bằng giữa
các người dùng. Một thuật toán phổ biến là Token Bucket: mỗi client được cấp một "thùng"
token với sức chứa tối đa, token được bổ sung đều đặn theo thời gian và mỗi request tiêu thụ
một token --- nhờ đó cho phép các đợt truy cập dồn (burst) ngắn hạn nhưng vẫn giới hạn được
tốc độ trung bình về dài hạn. Cơ chế này đặc biệt hữu ích để bảo vệ các endpoint nhạy cảm
như đặt hàng hay mua flash sale.

\subsection{Khả năng quan sát --- ba trụ cột giám sát hệ thống phân tán}

Trong một hệ thống gồm nhiều service, khi một request thất bại thì việc xác định lỗi phát
sinh từ đâu là rất khó nếu thiếu công cụ quan sát phù hợp. Khả năng quan sát (observability)
thường được xây dựng trên ba trụ cột bổ sung cho nhau:

\textbf{Metrics}: thu thập định kỳ các giá trị số phản ánh tình trạng hệ thống --- như tốc
độ request, độ trễ theo phân vị (p50/p95/p99), tỷ lệ lỗi hay mức tiêu thụ bộ nhớ. Metrics
phù hợp cho việc cảnh báo theo ngưỡng và lập kế hoạch năng lực, nhưng không giải thích được
nguyên nhân của từng sự kiện riêng lẻ.

\textbf{Distributed Tracing}: gán một định danh truy vết (TraceId) duy nhất cho mỗi request
và truyền nó qua tất cả service mà request đi qua. Mỗi đoạn xử lý tại một service tạo thành
một Span với thời điểm bắt đầu, kết thúc và các nhãn mô tả. Tập hợp các Span của cùng một
Trace tái dựng chính xác hành trình của request xuyên suốt hệ thống, giúp xác định service
nào đang gây ra độ trễ. Việc truyền ngữ cảnh truy vết giữa các service thường tuân theo chuẩn
W3C Trace Context.

\textbf{Log tương quan theo ngữ cảnh truy vết}: toàn bộ service sử dụng chung một định dạng
log thống nhất, trong đó mỗi dòng log được nhúng kèm tên service và định danh truy vết
(traceId). Nhờ traceId xuất hiện đồng nhất trên mọi dòng log của cùng một request, ta có thể
lọc và ghép toàn bộ log liên quan đến một request xuyên suốt các service, đồng thời đối chiếu
trực tiếp giữa log và trace tương ứng khi điều tra sự cố.

\textbf{Scheduled Jobs và Distributed Lock.} Một số tác vụ trong hệ thống phải chạy định kỳ
theo thời gian, chẳng hạn quét bảng outbox để phát sự kiện, hủy các đơn hàng quá hạn thanh
toán hay đối chiếu số liệu flash sale. Khi một service chạy nhiều instance song song, cần
bảo đảm chỉ một instance thực thi tác vụ định kỳ tại mỗi thời điểm để tránh xử lý trùng. Một
giải pháp phổ biến là Distributed Lock dựa trên Redis: cơ chế khóa bảo đảm chỉ một instance
giành được quyền thực thi, đồng thời khóa tự động được giải phóng sau một thời hạn nếu
instance đang giữ khóa gặp sự cố, nhờ đó tránh được tình trạng khóa kẹt vĩnh viễn.
\end{document}
